\section{GNU profiler (gprof)} 
\label{section:localisation:gprof}

\index{Profiler}
\index{Profiler!Instrumentation}
Every Linux system is shipped with the GNU profiler (\texttt{gprof}) which is
therefore a convenient starting point for every a performance analysis.
\texttt{gprof} is based upon instrumentation:
You ask the C/C++ or Fortran compiler to instrument
(Section~\ref{section:terminology:hpc-terms}) your code by passing in
\texttt{-p} to both the compiler and the linker step.
After that, the code's execution will monitor the behaviour and dump the
observation into a file \texttt{gmon.out}.
By typing in \texttt{gprof ./code}, you tell the profiler to take the
executable's debug information (remember to translate with \texttt{-g}) and
bring it together with the profile to assemble some meaningful stats.


The outcome consists of two major sections: 
The code first enlists the hotspots, i.e.~the most expensive functions, and then
it provides a call graph.
All the output is purely text-based by default.
There is no need for a fancy GUI or some X-forwarding.



\paragraph*{Overview}
 
\begin{itemize}[noitemsep]
  \item[$\boxplus$] Available everywhere.
  \item[$\boxplus$] All text-based.
\end{itemize}
 
 
\begin{itemize}[noitemsep]
  \item[$\boxminus$] You have to recompile.
  \item[$\boxminus$] Limited filter features to reduce the measurement overhead.
  \item[$\boxminus$] Cannot distinguish between different threads, and is not so
  straightforward to use with MPI. 
  \item[$\boxminus$] Focuses on timings and call counts only.
\end{itemize}

